\chapter{Ethical Issues}
\index{Ethical Issues}
\label{sec:Ethical Issues}

\section{Overview}
Geriatric care raises complex ethical challenges:

\textbf{Informed consent:} Must be obtained from patients with capacity; for those lacking capacity, consent is obtained from the legally authorized representative.

\begin{itemize}[noitemsep]
  \item \textbf{Decision-making capacity:} Capacity is task-specific and can be assessed using the Aid to Capacity Evaluation (ACE) tool. A patient with capacity can understand relevant information, appreciate the situation and its consequences, reason about treatment options, and communicate a choice. Incapacity requires surrogate decision-making based on substituted judgment or best interests.
  \item \textbf{Truth-telling:} Geriatric patients should be informed of their diagnosis and outlook in a sensitive, individualized manner. Withholding information (therapeutic privilege) is rarely justified.
  \item \textbf{Euthanasia and assisted dying:} Medical assistance in dying (MAiD) is legal in several jurisdictions. In Canada, MAiD is available to adults with a grievous and irremediable medical condition causing enduring suffering, with a reasonably foreseeable natural death. Clinicians have an ethical and legal duty to report suspected abuse to adult protective services.
\end{itemize}
\section{Causes}
\section{Risk Factors}

\begin{itemize}[noitemsep]
  \item Genetic predisposition and family history
  \item Advancing age
  \item Lifestyle factors (diet, exercise, smoking, alcohol)
  \item Environmental exposures
  \item Underlying medical conditions
  \item Medication use
\end{itemize}
\section{Signs and Symptoms}

\subsection{Common symptoms}
\begin{itemize}[noitemsep]
  \item \item Pain or discomfort in the affected area    \item Pain or discomfort in the affected area
  \end{itemize}

\subsection{Emergency warning signs}
\begin{itemize}[noitemsep]
  \item Severe or worsening symptoms of ethical issues require immediate medical evaluation
\end{itemize}

\section{Progression and Stages}
Advance requests for MAiD in dementia require careful ethical and legal consideration.
\begin{itemize}[noitemsep]
  \item \textbf{Elder abuse and neglect:} Affects 10--15\% of older adults, including physical, emotional, financial abuse, and neglect. The condition may progress through different stages of severity over time. Early stages often involve mild or subtle symptoms that may be overlooked. Moderate stages involve more noticeable symptoms affecting daily function. Advanced stages may involve significant impairment and complications.
\end{itemize}

\section{Complications}
\begin{itemize}[noitemsep]
  \item Disease progression leading to worsening symptoms
  \item Secondary infections or injuries
  \item Side effects of treatment
  \item Reduced quality of life
  \item Emotional and psychological impact
\end{itemize}
\section{Diagnosis}
Comprehensive medical history and physical examination Laboratory tests to assess organ function and rule out other conditions Imaging studies to visualize affected structures Specialized testing depending on the affected system Consultation with relevant specialists Monitoring of disease progression over time

\begin{itemize}[noitemsep]
  \item Comprehensive medical history and physical examination
  \item Laboratory tests to assess organ function
  \item Imaging studies to visualize affected structures
  \item Specialized testing depending on the affected system
  \item Consultation with relevant specialists
\end{itemize}
\section{Prevention}
\begin{itemize}[noitemsep]
  \item Maintain a healthy lifestyle with balanced diet and exercise
  \item Avoid known risk factors
  \item Attend regular medical check-ups for early detection
  \item Follow recommended screening guidelines
\end{itemize}
\section{Treatment Options}
Medications to manage symptoms and address underlying mechanisms Lifestyle modifications and self-care measures Physical therapy and rehabilitation Surgical intervention when indicated Regular monitoring and follow-up care Patient education and support services

\begin{itemize}[noitemsep]
  \item Medications to manage symptoms and address underlying mechanisms
  \item Lifestyle modifications and self-care measures
  \item Physical therapy and rehabilitation
  \item Surgical intervention when indicated
  \item Regular monitoring and follow-up care
\end{itemize}
\section{Living with It}
\begin{itemize}[noitemsep]
  \item Follow your treatment plan as prescribed
  \item Maintain regular follow-up with your healthcare provider
  \item Make necessary lifestyle adjustments
  \item Seek support from family, friends, or support groups
  \item Stay informed about your condition
\end{itemize}
\section{Outlook}
The outlook varies depending on the specific condition, severity at diagnosis, and response to treatment. Early intervention generally leads to better outcomes. Many conditions can be effectively managed with appropriate treatment. Regular follow-up care is important for monitoring and treatment adjustment.
